% Self-contained insertion after the proof of Theorem 1. Requires existing A_n,B_n.
\subsection{一般维数下公式族的渐近}
记
\[
 I_n(t)=\int_1^t\left(1+\frac1{2t}-\frac1{2\rho}\right)^{n-1}
                 \frac{\dint\rho}{\rho},\qquad
 F_n(t)=\bigl(\sqrt{A_n(t)}+\sqrt{B_n(t)}\bigr)^2,
 \qquad U_n=\inf_{t>1}F_n(t).
\]
这里 $U_n$ 是本文上界公式族的下确界；不预先把它与最优算子常数 $C_n$ 等同。
\begin{proposition}\label{prop:family-asymptotic}
当 $n\to\infty$ 时，
\[
 U_n=n\log n-n\log\log n+
 O\left(\frac{n\log\log n}{\log n}\right).
\]
特别地，$\limsup_{n\to\infty}C_n/(n\log n)\le1$。
\end{proposition}
\begin{proof}
先证明对全部 $n\ge2$、$t>1$ 一致成立的估计
\begin{equation}\label{eq:uniform-In}
 I_n(t)\le\log_+(t/n)+8+\log2,
 \qquad \log_+s=\max(0,\log s).
\end{equation}
在 $1\le\rho\le t/2$ 上，记
$\ell_t(\rho)=1+1/(2t)-1/(2\rho)$，则
\[
 \ell_t(\rho)^{n-1}\le
 e^{-(n-1)/(4\rho)}\le e^{-n/(8\rho)}.
\]
该区间中 $\rho\le n$ 的部分积分不超过
$\int_0^n(8/n)\,\dint\rho=8$，其中使用 $e^{-s}\le1/s$；
$\rho\ge n$ 的部分不超过 $\log_+(t/n)$。
余下区间 $[\max(1,t/2),t]$ 的对数长度不超过 $\log2$，
且 $0<\ell_t\le1$，故得到\eqref{eq:uniform-In}。
空区间在这里不贡献积分，因此该估计也覆盖 $t$ 靠近 $1$ 的情形。

令 $K=(8+\log2)/2$、$c=\log2$、$M=4/c$。
我们证明对每个 $n\ge16$，
\begin{equation}\label{eq:explicit-family-asymptotic}
 \begin{split}
 n\log n-n\log\log n-(K+\log M)n
 &\le U_n\\
 &\le n\log n-n\log\log n+(4+\log2)n.
 \end{split}
\end{equation}
为证上界，取 $t=2n/\log n$。于是 $A_n(t)\le\sqrt n$，
而
\[
 B_n(t)\le1+n\log t
 =1+n\log n-n\log\log n+n\log2
 \le1+n\log n.
\]
由 $\log n\le\sqrt n$，有
$\sqrt{A_nB_n}\le\sqrt{n^2+\sqrt n}\le n+1$。
展开 $F_n=A_n+B_n+2\sqrt{A_nB_n}$ 即得所需上界。

为证下界，由 $1-((1+1/t)/2)^n\ge0$ 和
\eqref{eq:uniform-In}，
\[
 B_n(t)\ge n\log t-\frac n2\log_+(t/n)-Kn.
\]
若 $t\ge n$，右端不低于 $n\log n-Kn$。
若 $1<t\le n$，置 $x=n/t\in[1,n)$。
凹性给出 $\log(1+s)\ge cs$（$0\le s\le1$），故
\[
 F_n(t)\ge A_n(t)+B_n(t)
 \ge n\log n+e^{cx}-n\log x-Kn.
\]
若 $x\le M\log n$，则
$\log x\le\log\log n+\log M$；若 $x>M\log n$，则
$e^{cx}>n^4\ge n\log n\ge n\log x$。
这两种情况均给出\eqref{eq:explicit-family-asymptotic}左侧。
下面收紧余项。由 $\ell_t(\rho)^{n-1}\le
\exp[-(n-1)(1/\rho-1/t)/2]$，置 $v=1/\rho-1/t$，得
\[
 I_n(t)\le\int_0^{1-1/t}
 \frac{e^{-(n-1)v/2}}{v+1/t}\,\dint v
 \le\frac{2t}{n-1}.
\]
令 $L=\log n$、$T=nL-n\log L$，并只考虑 $L\ge32$。
由式\eqref{eq:B-derivative}，$B_n$ 严格增加。
置 $x=n/t$、$x_0=L+4\log L\le2L$。
当 $x\le L/2$ 时，
$B_n(t)\ge B_n(2n/L)\ge T+n\log2-4n/L\ge T$。
当 $L/2\le x\le x_0$ 时，
\[
 B_n(t)\ge nL-n\log x-2n/x
 \ge T-\frac{4n(\log L+1)}L.
\]
当 $x\ge x_0$ 时，由 $x_0^2\le4L^2\le2e^L$ 及
$\log(1+s)\ge s-s^2/2$，
\[
 A_n(t)\ge(1+x_0/n)^n\ge e^{x_0-x_0^2/(2n)}
 \ge nL^4/e\ge nL\ge T.
\]
这就证明 $U_n\ge T-4n(\log L+1)/L$。
反过来，取 $t=n/(L-3\log L)$。
因为 $L-3\log L\ge L/2$，有
\[
 A_n(t)\le n/L^3,\qquad
 B_n(t)\le T+6n\log L/L+1,\qquad B_n(t)\le2nL.
\]
于是 $F_n(t)\le T+n(6\log L+5)/L$，得到所述强化余项。
最后由 $C_n\le U_n$ 得到所述上极限不等式。
注意这里的下界仅针对 $U_n$，不构成 $C_n$ 的下界。
\end{proof}

一维情形须单独说明：$I_1(t)=\log t$，且
\[
 A_1(t)=1+1/t,\qquad B_1(t)=\tfrac12(1-1/t+\log t).
\]
因 $\log t>1-1/t$（$t>1$），有 $F_1(t)\ge A_1(t)+B_1(t)>2$；
而 $t\downarrow1$ 时 $F_1(t)\to2$。
所以 $U_1=2$，但此下确界不在 $t>1$ 内达到。

\subsection{参数导数与有限膨胀参数的比较}
对 $n\ge2$，Leibniz 法则给出
\begin{equation}\label{eq:B-derivative}
 B_n'(t)=\frac n{2t}
 +\frac n{2t^2}\left(\frac{1+1/t}{2}\right)^{n-1}
 +\frac{n(n-1)}{4t^2}\int_1^t
       \ell_t(\rho)^{n-2}\frac{\dint\rho}{\rho}>0.
\end{equation}
当 $n=1$ 时，单独有 $B_1'(t)=1/(2t)+1/(2t^2)$。
由此
\begin{equation}\label{eq:F-derivative}
 F_n'(t)=(\sqrt{A_n}+\sqrt{B_n})
 \left[\frac{B_n'(t)}{\sqrt{B_n(t)}}
 -\frac n{t^2}(1+1/t)^{n/2-1}\right].
\end{equation}
式\eqref{eq:F-derivative}给出内点极小值的必要条件，
并不证明全局极小点或驻点的唯一性。
事实上 $B_n(t)=n(t-1)+O((t-1)^2)$，所以 $F_n'(t)\to+\infty$
当 $t\downarrow1$。
对固定 $n$，$B_n(t)\ge1-((1+1/t)/2)^n+(n/2)\log t\to\infty$；
因而若某点满足 $F_n(t)<2^n$，则下确界必在内点取得。

最后说明，在当前预算估计中保留有限 $a$ 不会改善固定 $t$ 的上界。
置 $u=1/a$、$L=\log t$，并令
\[
 A_{n,a}=(1+1/t+u)^n,\qquad E_{n,a}=a^{-n}+\alpha^{-n}J_a(t).
\]
利用 $\ell_a=\alpha\ell_t$ 可得精确恒等式
\[
 E_{n,a}-B_n
 =u^n+[(1+u)^n-1](1+nL)-\frac n2uI_n(t).
\]
由 $I_n(t)\le L$ 及 $(1+u)^n-1\ge nu$，
\[
 E_{n,a}-B_n\ge u^n+nu[1+(n-\tfrac12)L]>0.
\]
同时 $A_{n,a}>A_n$，故每个有限 $a>2t$ 的 $\tau$ 优化值
$(\sqrt{A_{n,a}}+\sqrt{E_{n,a}})^2$ 严格大于 $F_n(t)$。
这只排除了在本预算族中靠有限 $a$ 改善的路线；
联合下确界仍等于 $U_n$，因为固定 $t$ 时可令 $a\to\infty$。
